\documentclass[11pt]{article}
\usepackage[margin=25mm]{geometry}
\usepackage{amsmath,amssymb,amsthm,booktabs}
\usepackage[colorlinks=true,urlcolor=blue,linkcolor=blue]{hyperref}
\usepackage[T1]{fontenc}
\newtheorem{theorem}{Theorem}
\newtheorem{proposition}{Proposition}
\setlength{\parindent}{0pt}
\setlength{\parskip}{6pt}
\setlength{\emergencystretch}{3em}
\begin{document}
\begin{center}
{\Large A single edge density does not determine graph statistics}\\[5pt]
{\large Conjecture 00000004398}\\[5pt]
zbw123456 \quad 14 October 2026
\end{center}
\textbf{Scope.} We use the printed definition: a forcing graph has a single
density determining all graph statistics. Density means homomorphism density.
This note refutes the proposed inclusion of all complete graphs under that
reading. The Lean theorem concerns finite simple graphs; it does not assert a
formalization of graphon integration or asymptotic quasirandomness.

For finite simple graphs $F,G$ with $|V(G)|>0$, put
\[
 t(F,G)=\frac{\#\{f:V(F)\to V(G): f\text{ preserves edges}\}}
 {|V(G)|^{|V(F)|}}.
\]
A necessary consequence of the proposed classification is that equality of
$t(K_2,G)$ and $t(K_2,H)$ implies equality of $t(J,G)$ and $t(J,H)$ for every
finite simple graph $J$.

\begin{theorem}
The graphs $G=K_3\sqcup K_1$ and $H=K_{1,3}$ have the same $K_2$ density
and different $K_3$ densities. Consequently $K_2$ does not have the stated
forcing property, and the proposed classification is false under that definition.
\end{theorem}
\begin{proof}
Both graphs have four vertices and three edges. Every undirected edge gives
exactly two homomorphisms from $K_2$, so both homomorphism counts are six.
A homomorphism from $K_3$ to a loopless graph must be injective. The triangle
of $G$ gives $3!=6$ homomorphisms; the star $H$ has no triangle and gives none.
Thus
\[
 t(K_2,G)=t(K_2,H)=\frac6{16}=\frac38,
 \qquad t(K_3,G)=\frac6{64}=\frac3{32}\ne0=t(K_3,H).
\]
Since $K_2$ is complete, it contradicts the forward implication of the
claimed classification. No classification of the remaining graphs is needed.
\end{proof}

\textbf{Formal correspondence.} The file \texttt{GraphDensity.lean} defines
Boolean adjacency on \texttt{Fin n}, simplicity, complete graphs, all vertex
maps, the edge-preservation test and homomorphism counts. Density equality
uses cross multiplication of positive integer denominators.
\texttt{mem\_allMaps} proves coverage of every vertex map;
\texttt{enumeration\_two} and \texttt{enumeration\_three} verify sizes 16 and 64
and absence of duplicate tuples in the exact enumerations used here.
The final theorem \texttt{complete\_two\_not\_forcing} applies the purported
property to these actual graphs; \texttt{not\_all\_complete\_forcing} refutes
the necessary universal implication. Only the standard axiom \texttt{propext}
appears in the final theorem audit.

\textbf{Interpretation limit.} The usual quasirandom-forcing terminology is
more specific than the problem's short definition. This note does not silently
replace it with a new asymptotic assertion. The finite computation is exact;
acceptance of the printed-definition reading remains for the competition reviewers.

\textbf{Sources.} Official bilingual conjecture is included as
\texttt{conjecture.md}. For the conventional homomorphism-density definition,
see Grzesik, Kr\'al' and Pikhurko,
\href{https://arxiv.org/abs/2303.04041}{Forcing Generalized Quasirandom Graphs Efficiently}.
\end{document}
